%% ============================================================
%%  Performance Testing and Predictors
%%  Merged from physical_testing.tex (§A, §B, §E, §F) and
%%  performance_predictors.tex (§C, §D)
%% ============================================================

\section{Performance Testing and Predictors}
\label{sec:predictors}

This section synthesises evidence on field measurement methodology (Section~\ref{sec:ptestmeasure}), finger strength as the dominant grade predictor (Section~\ref{sec:mfspredictor}), aerobic capacity and critical force (Section~\ref{sec:cfpredictor}), anthropometric and other physical predictors (Section~\ref{sec:anthropometrics}), multivariate models and probabilistic grade estimation (Section~\ref{sec:multivariate}), and test batteries used in research and applied practice (Section~\ref{sec:testbatteries}). All correlations reported are from single-time-point cross-sectional designs unless explicitly noted; no causal inference is warranted without controlled longitudinal or intervention data.

%---------------------------------------------------------------------
\subsection{Physical Measurement Methodology}
\label{sec:ptestmeasure}

\subsubsection{Maximum Finger Strength (MFS): Measurement Methods}

Reliable MFS assessment requires standardisation of edge depth, body position, grip type, contraction protocol, and the force-measurement device. A systematic review of 15 studies \cite{TEMP_Jullien2025} found ICC values ranging from 0.294 to 1.0 with a median of 0.86; studies using 20--23~mm edges consistently achieved ICC $>$ 0.90, establishing this edge depth as the recommended research standard. Edges outside this range---the literature spans 6--60~mm across protocols---introduce large inter-study heterogeneity that precludes direct comparison of MFS values.

\medskip
\noindent\textbf{Force-measurement device.} Torr et al.\ \cite{TEMP_Torr2022} validated a consumer linear-position transducer (Tindeq Progressor) against a Kistler force plate ($n = 229$), reporting ICC = 0.99 for one-handed rung pull and bent-arm lock-off. Within-day reliability ICC = 0.90--0.98; between-day ICC = 0.95--0.98. Tindeq-based MFS measurement is therefore suitable for both applied and research settings. Tindeq RFD measurement is addressed separately below.

\medskip
\noindent\textbf{Body position.} Stien et al.\ \cite{TEMP_Stien2025pos} compared five body positions and reported that the seated bent-arm position at 90$^{\circ}$ elbow flexion (Sit90) produced the strongest correlation with climbing performance: bouldering $r = 0.73$, sport-climbing $r = 0.64$.\singlesource\ The straight-arm hanging position yielded higher $r^{2}$ with grade in some analyses but lower test reliability (ICC = 0.88 straight-arm vs.\ 0.94 bent-arm), making Sit90 the preferred protocol for serial testing. Because this study is currently a single-source DOI (unverified), the position-difference finding carries \toverify status pending independent replication.

\medskip
\noindent\textbf{Grip type.} The half-crimp (PIP and DIP joints at $\sim90^{\circ}$ and $\sim45^{\circ}$) is the recommended research standard: it produces higher absolute force than open-hand on identical edges, is ecologically valid for most climbing contexts, and permits direct comparison across studies \cite{TEMP_Jullien2025}. Open-hand (slope) grip generates lower A2 pulley stress and is therefore argued as safer for repeated testing sessions. One \toverify source\singlesource\ (Wolf et al.\ 2025, single source) reports that on a 23~mm rung, open-hand produced a mean 9.8\% higher force than half-crimp in advanced climbers and advocates open-hand to reduce pulley loading; this finding must not be treated as settled guidance until independently replicated.

\medskip
\noindent\textbf{Minimal Detectable Change.} No peer-reviewed MDC value for climbing MFS has been published. From the ICC range reported by Jullien et al.\ \cite{TEMP_Jullien2025} (ICC 0.86 median; 0.90 for 20--23~mm edges), the estimated MDC is approximately 5--7\% for within-session repeated testing under standardised conditions. This estimate should be applied cautiously; a formal SEM-derived MDC from a large climbing-specific sample remains a research gap.

\subsubsection{MFS Testing Parameters: Summary}

{\footnotesize
\setlength{\tabcolsep}{5pt}
\renewcommand{\arraystretch}{1.2}
\begin{tabular}{>{\raggedright\arraybackslash}p{3.2cm}
                >{\raggedright\arraybackslash}p{3.5cm}
                >{\raggedright\arraybackslash}p{3.5cm}
                >{\raggedright\arraybackslash}p{4.0cm}}
\toprule
\textbf{Parameter} & \textbf{Recommended standard} & \textbf{ICC (reported range)} & \textbf{Notes} \\
\midrule
Edge depth & 20--23~mm & $>$0.90 at 20--23~mm & Edges $<$18~mm or $>$25~mm show lower ICC \cite{TEMP_Jullien2025} \\[3pt]
Body position & Sit90 (seated bent-arm, 90$^{\circ}$ elbow) & 0.94 & Highest performance correlation $r = 0.73$ bouldering \cite{TEMP_Stien2025pos} (\toverify) \\[3pt]
Grip type & Half-crimp & --- & Ecological validity; higher absolute force than open-hand; open-hand lower pulley stress \\[3pt]
Device & Tindeq Progressor or equivalent force plate & 0.99 vs.\ Kistler \cite{TEMP_Torr2022} & 80~Hz sampling; not adequate for RFD (see below) \\[3pt]
Contraction protocol & 3 maximal efforts, 90~s rest between, best or mean recorded & --- & $\geq$2 familiarisation sets recommended \\[3pt]
MDC estimate & $\sim$5--7\% & --- & No published climbing-specific MDC; formal derivation needed \\
\bottomrule
\end{tabular}
}

\begin{evidencebox}{MFS Testing: Standard Protocol · EVIDENCE\_BACKED · ESS 7}
Based on Jullien et al.\ \cite{TEMP_Jullien2025} (15-study reliability review, ICC median 0.86, ICC $>$ 0.90 for 20--23~mm) and Stien et al.\ \cite{TEMP_Stien2025pos} (position comparison, \toverify): measure MFS as the peak force from 3 maximal isometric efforts on a 20--23~mm edge, half-crimp grip, Sit90 body position (seated, elbow 90$^{\circ}$), with 90~s rest between trials. Record best value and mean. Express as absolute force (N or kg) and relative to body mass (\%BW). A change $\geq$5--7\% is required to exceed estimated measurement noise. Tindeq Progressor is a validated device for MFS (ICC = 0.99) \cite{TEMP_Torr2022} but is not adequate for RFD research due to 80~Hz sampling limitation.
\end{evidencebox}

\subsubsection{Rate of Force Development (RFD) Measurement}

RFD discriminates elite from advanced climbers in a way that peak MFS alone does not. Stien et al.\ \cite{SVS+21} reported effect sizes of ES = 1.02--1.58 between elite and advanced/intermediate climbers at relative time periods (10--90\% of time-to-peak-force), whereas no significant RFD difference was detected between advanced and intermediate groups. This suggests RFD becomes a meaningful discriminator only at the elite level.

\medskip
\noindent\textbf{Time-window convention.} Absolute millisecond windows (e.g., RFD$_{100}$, RFD$_{200}$) conflate differences in time-to-peak-force with differences in force production rate. Relative time-period analysis (RFD at 10\%, 30\%, 50\%, 70\%, 90\% of contraction time) is recommended \cite{SVS+21}. Vereide et al.\ \cite{VAH+22} confirmed that both peak force and RFD differ significantly by performance level when measured with adequate temporal resolution.

\medskip
\noindent\textbf{Sampling-rate requirement.}

\begin{bioplausbox}{Tindeq 80~Hz is Insufficient for RFD Research · BIOPLAUSIBLE · ESS 3}
The Tindeq Progressor samples at 80~Hz (12.5~ms resolution), which is insufficient to resolve the early ascending phase of a maximal isometric finger-strength contraction (peak force typically reached in 300--600~ms). Estimated MDC for Tindeq-derived RFD is 49--76\%, rendering most within-subject RFD comparisons statistically uninformative at this sampling rate. Research-grade RFD quantification requires $\geq$200~Hz sampling. Applied practitioners tracking serial RFD with Tindeq should treat trend directions only, not absolute values; no peer-reviewed study has validated Tindeq specifically for climbing RFD (classification: \bioplausible, mechanistic extrapolation from general force-platform sampling theory; no climbing-specific RFD validation study at 80~Hz confirmed).
\end{bioplausbox}

%---------------------------------------------------------------------
\subsection{Finger Strength as Predictor of Climbing Grade}
\label{sec:mfspredictor}

\subsubsection{Core Predictor Evidence}

MFS assessed on a climbing-specific edge is the strongest single measurable predictor of rock-climbing performance across multiple study designs and populations. Buraas and Stien \cite{BBW25} ($n = 19$ males, 6b+--8c Font, observational cross-sectional) reported climbing-specific finger strength correlations of $r = 0.89$ with bouldering grade, $r = 0.67$ with redpoint grade, and $r = 0.82$ with campus board performance, compared with only $r = 0.54$ for a standard handgrip dynamometer on redpoint grade---confirming that climbing-specific testing outperforms general grip dynamometry for predicting climbing ability.

In the largest normative sample to date, Stien et al.\ \cite{BMK+25} ($n = 307$, 185~M + 122~F; cross-sectional) reported that finger strength explained 68\% of climbing ability variance in males and 64\% in females in a univariate model; intermittent endurance added a further 28\% and 34\%, respectively. The commonly used finger hang test ($R \approx 0.75$ with grade) was found to be 65--80\% strength-dependent, indicating that endurance tests with heavy strength loading confound pure endurance measurement.

Torr et al.\ \cite{TEMP_Torr2022} ($n = 229$, mixed sex, cross-sectional) reported relative MFS correlations of $r = 0.42$--$0.50$ with climbing performance across grip types and body positions, broadly consistent with larger cohorts.

\subsubsection{Relative vs.\ Absolute MFS}

\begin{evidencebox}{Relative MFS as Dominant Grade Predictor · EVIDENCE\_BACKED · ESS 7}
Relative MFS (\% body weight) is universally a stronger predictor of climbing grade than absolute MFS (N or kg) across all reviewed studies. The mechanistic rationale is direct: climbing requires moving body mass against gravity, so performance scales with force-to-weight ratio rather than raw force output. Stien et al.\ \cite{TEMP_Stien2025pos} report relative MFS $r = 0.425$--$0.767$ with climbing grade vs.\ absolute MFS $r = 0.422$--$0.741$ depending on body position and grade metric; the relative advantage is consistent if not always large in magnitude. Express all training targets and normative benchmarks as \%BW.
\end{evidencebox}

\medskip
\noindent\textbf{Grip technique and edge depth.} Half-crimp on a 20~mm edge (\toverify; de facto research and applied standard via Lattice Training and academic protocols) yields $R^{2} = 0.48$--$0.58$ for bouldering in a 6-technique comparison \cite{BMK+25}, the highest of any single technique tested. MacKenzie et al.\ \cite{TEMP_MacKenzie2020} and Grant et al.\ \cite{GHWA96} corroborate finger strength as a key discriminator between performance levels in multi-variable models.

\subsubsection{Multivariate Performance Models}

Table~\ref{tab:multivariate} summarises multivariate models predicting climbing grade from physical testing data.

{\footnotesize
\setlength{\tabcolsep}{5pt}
\renewcommand{\arraystretch}{1.2}
\begin{longtable}{>{\raggedright\arraybackslash}p{3.2cm}
                  >{\raggedright\arraybackslash}p{1.0cm}
                  >{\raggedright\arraybackslash}p{2.0cm}
                  >{\raggedright\arraybackslash}p{8.0cm}}
\caption{Multivariate models predicting climbing performance from physical testing.}\label{tab:multivariate}\\
\toprule
\textbf{Study} & \textbf{$n$} & \textbf{$R^{2}$ total} & \textbf{Predictors and relative contribution} \\
\midrule
\endfirsthead
\multicolumn{4}{c}{\textit{(continued)}}\\
\toprule
\textbf{Study} & \textbf{$n$} & \textbf{$R^{2}$ total} & \textbf{Predictors and relative contribution} \\
\midrule
\endhead
\bottomrule
\endlastfoot

Mermier et al.\ \cite{MJPS00} & 44 & 0.66 & Training volume 58.9\%; flexibility 19.8\%; span-to-height ratio 9.8\%; remaining 11.5\% unexplained. Finger strength not separately isolated; training volume proxies accumulated finger loading \\[3pt]

Bala\v{s} et al.\ \cite{BPMC11} & ${\sim}$60 & ${\sim}$0.97$^{a}$ & Hand-arm strength + finger endurance primary predictors; model $R^{2}$ high but SEM-based — see note $a$ \\[3pt]

Stien et al.\ \cite{BMK+25} & 307 & ${\sim}$0.96 & Finger strength 68\% (M) / 64\% (F); intermittent endurance adds 28\% (M) / 34\% (F); model validated on normative sample \\[3pt]

\end{longtable}
\noindent$^{a}$\textit{Bala\v{s} et al.\ \cite{BPMC11}: the reported $R^{2} \approx 0.97$ was derived via stepwise regression on a small sample ($n \approx 60$); overfitting is likely and the model has not been externally cross-validated. Treat as exploratory.}
}

\subsubsection{Cross-Sectional Design Limitations}

\begin{evidencebox}{Finger Strength as Primary Predictor: Limitations · EVIDENCE\_BACKED · ESS 5}
All MFS--grade correlation studies reviewed are cross-sectional \cite{MJPS00, BPMC11, BMK+25, BBW25, TEMP_Torr2022}. Three constraints follow: (1) reverse causation is plausible --- climbers with higher grades have trained more, accumulating greater finger strength; the correlation does not isolate the causal contribution of MFS to grade attainment. (2) Climbing grades are self-reported in most studies, introducing recall and sandbag bias. (3) All samples are convenience cohorts (volunteer, club, or gym-recruited), limiting generalisability to casual or elite competition populations. Intervention evidence (hangboard RCTs, Section~\ref{sec:maxvsrep}) provides causal support for MFS training improving performance, but no RCT has used on-wall grade as a primary outcome.
\end{evidencebox}

\subsubsection{Community Benchmarks}

\begin{folklorebox}{Lattice Normative MFS Benchmarks (7-s 2-Arm Hang, 20~mm, Half-Crimp, \%BW) · FOLKLORE\_VERIFIED}
The following benchmarks originate from Lattice Training applied data and are reproduced across $\geq$3 independent climbing-community sources (Lattice Training blog, WMG Climbing, UKC forums). Values represent 7-second 2-arm dead-hang on a 20~mm edge at \%BW:

\smallskip
\begin{tabular}{lll}
\toprule
\textbf{Bouldering grade} & \textbf{Font equiv.} & \textbf{\%BW (2-arm, 7~s, 20~mm)} \\
\midrule
V4  & 6b+ & 49\% \\
V6  & 7a  & 61\% \\
V8  & 7b+ & 73\% \\
V10 & 7c+ & 85\% \\
V12 & 8b  & 96\% \\
V14 & 8c  & 106\% \\
V17 & 9a  & 118\% \\
\bottomrule
\end{tabular}

\smallskip
These are 2-arm values suitable for V4--V12 profiling. From approximately V14 (8c) onwards, Lattice transitions to 1-arm profiling because bilateral 2-arm hang ceiling effects (values approaching or exceeding 100\%BW) compress the discrimination range at elite levels. Direct comparison between these benchmarks and individual test data requires identical protocol: same edge depth, same position, same contraction duration.
\end{folklorebox}

\subsubsection{Pull-Up Strength as Secondary Predictor}

Pull-up strength is a moderate predictor of bouldering performance: Buraas and Stien \cite{BBW25} report $r = 0.55$ for bouldering grade ($n = 19$), secondary to finger strength ($r = 0.89$). For sport climbing at the elite level, one \toverify source\singlesource\ (Lattice Training 2025 performance data blog post) reports that the 2-repetition-maximum pull-up was the strongest single predictor in a route-climbing model ($R = 0.521$, adj.\ $R^{2} = 0.253$), potentially reflecting the greater importance of sustained upper-body endurance on long routes. This finding is blog-level only and must not be treated as settled guidance; a peer-reviewed analysis of the same dataset has not been published. Pull-up training and testing protocols are addressed in detail in the physical test battery section.

\subsection{Aerobic Capacity, Critical Force, and Endurance as Predictors}
\label{sec:cfpredictor}
% -------------------------------------------------------------------

\subsubsection{Critical Force: Theoretical Framework}

Critical Force (CF) is operationally defined as the highest sustainable intermittent isometric force that can be maintained indefinitely without continuous depletion of anaerobic stores; it represents the isometric analogue of the critical power construct.
W$'$ (``W-prime'') is the finite anaerobic work capacity above CF, expressed in N$\cdot$s or kg$\cdot$s.
Together CF and W$'$ fully parameterise the two-parameter work-time model for intermittent isometric effort:

\[
  W_{\text{total}} = \mathrm{CF} \cdot t + W'
\]

where plotting total isometric work ($W_{\text{total}}$) against time to exhaustion ($t$) yields a straight line whose slope is CF and whose y-intercept is W$'$.

CF is conceptually distinct from maximal finger strength (MFS): CF defines the aerobic ceiling of intermittent contraction capacity (${\approx}40$--$42\,\%$ MVC in trained climbers), whereas MFS defines peak force capacity.
They describe orthogonal physiological capacities; a climber may have high MFS but a low CF:MFS ratio, indicating limited aerobic recovery between hard moves.
Giles et al.\ (2019) measured mean CF of $425.7 \pm 82.8$\,N ($= 41.0 \pm 6.2\,\%$ MVC) and W$' = 30{,}882 \pm 11{,}820$\,N$\cdot$s in a mixed climber cohort \cite{GCT+19}.

\subsubsection{CF Measurement Protocols}

\textbf{Gold-standard multi-trial protocol} \cite{GCT+19}: participants perform three to four separate sessions at fixed force targets (45\,\%, 60\,\%, and 80\,\% MVC) using a 7:3 work-to-rest ratio until failure; CF and W$'$ are derived from the linear regression of total work on time to exhaustion.
The protocol requires $\geq 3$ sessions separated by $\geq 48$\,h, limiting field applicability.

\textbf{All-out single-session protocol} \cite{GHM+21}: 24 maximal one-arm contractions at 7:3 work-to-rest ratio; CF is calculated as the mean force of the final 6 contractions; W$'$ is derived from total work above CF.
Test--retest reliability: ICC $= 0.900$ (95\,\% CI 0.616--0.979) for CF; ICC for W$'= 0.768$, indicating substantially lower repeatability for the anaerobic work-capacity parameter.
Bal\'{a}\v{s} et al.\ (2024) independently validated the all-out method, reporting ICC $> 0.90$ for CF \cite{BGJ+24}.

\textbf{Critical overestimation finding (TEMP\_Balas2024)}: CF as conventionally defined by the mean of the last three contractions overestimates the physiologically sustainable intensity.
In one validation study only 38\,\% of participants could sustain 720\,s of continuous loading at the calculated CF intensity; the \emph{end-force} (lowest recorded contraction value) performed better as a threshold for sustainable effort \cite{BGJ+24}.
Practitioners using CF-derived training zones should treat the standard calculation as an \emph{upper bound} on sustainable intensity, not a confirmed aerobic threshold.

\textbf{Tindeq all-out protocol}: Commercial Tindeq devices implement the Giles (2021) protocol algorithmically.
Community documentation is extensive (r/climbharder, Lattice Training blog), but no peer-reviewed study has validated the Tindeq-specific CF output against a criterion all-out test or the multi-trial gold standard.

\begin{bioplausbox}{Tindeq-Derived CF as a Training Metric · \bioplausible · ESS 2}
The Tindeq all-out CF output applies the Giles (2021) protocol to reliable force-time data; the underlying methodology has ICC $> 0.90$ for CF \cite{GHM+21, BGJ+24}.
Device-specific systematic bias against a criterion multi-trial protocol has not been published.
CF zones derived from Tindeq output should be interpreted with reference to the known overestimation of sustainable intensity noted above.
\end{bioplausbox}

\subsubsection{CF as a Route-Climbing Predictor}

\begin{evidencebox}{Critical Force as Route-Climbing Predictor · \evidencebacked · ESS 6}
Giles et al.\ (2019, $n$ not fully reported, mixed sex and grade): CF expressed as \%BW correlated $r = 0.61$ with sport climbing (route) grade; W$'$ correlated $r = 0.34$ with bouldering grade \cite{GCT+19}.
The stronger route association is mechanistically consistent with CF reflecting aerobic endurance capacity, which is demanded more by sustained route climbing than by maximal-effort bouldering problems.

Lattice Training (2025 internal dataset, $N = 422$ routes): CF entered a multiple regression as $R = 0.401$, adjusted $R^2 = 0.16$ for the full sample.
CF was \emph{not} a significant predictor at the Higher Elite sub-group level ($p = 0.052$), indicating diminishing predictive returns at top grades where all competitors have well-developed aerobic endurance.

Stien et al.\ (2025 normative, $n = 307$): intermittent endurance (CF-related measure) explained 28\,\% (males) and 34\,\% (females) of climbing grade variance \emph{after} controlling for finger strength \cite{BMK+25}, confirming that CF-related endurance and MFS are partially independent predictors.

\textbf{Limitations}: Giles (2019) sample characteristics and allocation details are not fully reported; grade self-report introduces noise; the Lattice dataset is unpublished and lacks peer review.
\end{evidencebox}

\textbf{CF:MFS ratio as a community heuristic}: The threshold of CF $> 50\,\%$ MFS for route climbers and $> 70\,\%$ for endurance specialists is widely cited in the climbing coaching community (Lattice Training blog; StrengthClimbing.com).
No peer-reviewed intervention has tested this threshold prospectively.

\begin{bioplausbox}{CF:MFS Ratio Thresholds for Route vs.\ Endurance Climbers · \bioplausible · ESS 2}
A CF:MFS ratio of $> 50\,\%$ (route) and $> 70\,\%$ (endurance) is mechanistically plausible as an indicator of aerobic recovery capacity relative to peak force, and is coherent with the observed differential predictiveness of CF vs.\ MFS across disciplines.
Threshold values originate from coaching practice, not RCT-derived cut-points.\singlesource
\end{bioplausbox}

\subsubsection{Grade-Specific Endurance Threshold}

\begin{folklorebox}{Grade-Band Endurance Thresholds · \folkloreverified · ESS 1}
Community consensus (Lattice Training blog; StrengthClimbing.com; r/climbharder aggregate discussion) holds that:
(1) below approximately V6\,/\,7b, technique dominates grade outcome and aerobic endurance is weakly predictive;
(2) between V8--V10\,/\,7c--8a, both finger strength and forearm aerobic endurance become critical for route success;
(3) above 8a$+$\,/\,Higher Elite, CF predictiveness diminishes (consistent with Lattice $p = 0.052$ finding above).
These boundaries represent coaching heuristics with partial empirical corroboration at the upper extreme; the lower thresholds lack quantitative validation.
\end{folklorebox}

\subsubsection{Systemic VO\textsubscript{2}max: Null Grade Predictor}

\begin{evidencebox}{Systemic VO\textsubscript{2}max: Null Grade Predictor · \evidencebacked · ESS 5}
Billat et al.\ (1995, $n$ not reported): ascent of 7b routes elicited only 37.7--45.6\,\% of treadmill VO\textsubscript{2}max, heart rate 77--85.5\,\% of maximum, and post-ascent blood lactate 4.3--5.7\,mmol/L \cite{TEMP_Billat1995}.
Multiple reviews confirm that treadmill- or cycle-ergometer-derived VO\textsubscript{2}max does not correlate significantly with climbing grade when measured in trained populations \cite{Wat04}.
The mechanistic basis is unambiguous: climbing performance is limited primarily by localised isometric finger flexor fatigue and forearm intramuscular oxygen delivery, not by whole-body cardiac or pulmonary capacity.
Training systemic aerobic fitness therefore cannot substitute for forearm-specific aerobic adaptation.
\end{evidencebox}

\subsubsection{Blood Lactate: Secondary Role}

Post-ascent venous lactate is consistently elevated (4--7\,mmol/L) after hard routes, reflecting local glycolytic contribution during crux sections.
No published study demonstrates that systemic lactate threshold independently predicts climbing grade when finger strength variables are controlled.
Local intramuscular lactate clearance and buffering capacity in the finger flexors are mechanistically relevant to inter-move recovery during sustained sequences, but the corresponding in-vivo measurement has not been validated as a grade predictor.

\begin{bioplausbox}{Local Forearm Lactate Clearance as Inter-Move Recovery Determinant · \bioplausible · ESS 2}
Mechanistic rationale: forearm intramuscular carnosine buffering and capillary lactate clearance determine how quickly the forearm can recover between sustained crux sequences.
Direct measurement of finger flexor--specific lactate kinetics as a grade predictor has not been performed; systemic lactate metrics are inadequate proxies for this local capacity.
\end{bioplausbox}

% -------------------------------------------------------------------
\subsection{Anthropometrics and Other Physical Predictors}
\label{sec:anthropometrics}
% -------------------------------------------------------------------

\subsubsection{Ape Index: Null Finding Within Competitive Samples}

\begin{evidencebox}{Ape Index: Null Finding Within Competitive Samples · \evidencebacked · ESS 5}
Ozimek et al.\ (2023) systematic review (18 studies): ``none of the analysed works showed differences between sport climber groups in the ape index'' \cite{TEMP_Ozimek2023}.
Mermier et al.\ (2000, $n = 44$, mixed sex): an anthropometric component (including ape index, height, body composition) explained only $0.3\,\%$ of climbing performance variance; a training-variable component (years climbing, training frequency, finger strength) explained $58.9\,\%$ \cite{MJPS00}.
Watts et al.\ (2003, $n = 90$ junior competitive climbers): no significant ape index--grade correlation was detected \cite{WJL+03}.
\end{evidencebox}

\begin{toverifybox}{Ape Index Advantage in International vs.\ National-Level Boulderers · \toverify · ESS 2}
Frontiers (2025, $n = 34$; IFSC international vs.\ national boulderers): ape index was significantly higher in the international group ($p < 0.05$) \cite{DTB+25}.
This finding directly contradicts the systematic review \cite{TEMP_Ozimek2023}.
Possible explanations include small sample size ($n = 34$), narrow grade band producing range restriction, ethnicity-related body proportions confounding the comparison, or genuine ape index advantage at the extreme upper end of competition level.
Interpret as preliminary pending replication in a larger, grade-matched sample.
\end{toverifybox}

\subsubsection{Body Composition}

Lower body fat percentage (\%BF) is associated with elite competitive status (EVIDENCE\_BACKED, ESS 5).
Elite male sport climbers typically present with 4--12\,\% BF by DEXA \cite{TEMP_Ozimek2023}; elite female climbers show correspondingly low values, though sex-specific reference ranges are less consistently reported.
However, Mermier et al.\ (2000) demonstrated that body mass explains $< 1\,\%$ of performance variance \emph{when training variables are entered simultaneously} \cite{MJPS00}, indicating that leanness is a characteristic of elite climbers rather than a causal bottleneck for non-elite populations.

Structural adaptations observed in elite climbers include A2 pulley hypertrophy and greater forearm flexor muscle belly thickness compared with recreational climbers.
The direction of causality is uncertain: heavy training volume produces both adaptations and elite-level performance simultaneously; neither has been established as a prerequisite.

\begin{toverifybox}{Weight Cutting as a Performance Strategy for Non-Elite Climbers · \toverify · ESS 0}
Weight reduction is not supported as a performance intervention for non-elite climbers by any published controlled trial.
The Sch\"{o}ffl research group has documented clinically significant rates of disordered eating and Relative Energy Deficiency in Sport (RED-S) in competitive junior and senior climbers \cite{TEMP_Schoffl_REDS}.\footnote{TEMP\_Schoffl\_REDS: specific Sch\"{o}ffl RED-S citation requires DOI verification --- add via \texttt{add\_citation} workflow before publication.}
Deliberate weight cutting in the absence of medical supervision carries documented harm risk (bone stress injury, hormonal disruption, performance impairment); it is classified \textbf{Very High} injury risk for non-elite populations.
\end{toverifybox}

\subsubsection{Summary of Other Physical Predictors}

Table~\ref{tab:physpredictors} summarises correlation data for additional physical variables commonly assessed in climber profiling.

\begin{table}[ht]
\centering
\caption{Physical predictor correlations with climbing grade. All data from cross-sectional designs; causal inference is not supported. ``Null'' indicates $r < 0.15$ or $p > 0.05$ in the largest available sample.}
\label{tab:physpredictors}
\small
\begin{tabular}{p{3.5cm} p{2.2cm} p{2.5cm} p{2.0cm} p{3.5cm}}
\hline
\textbf{Predictor} & \textbf{Correlation} & \textbf{Sample} & \textbf{ESS} & \textbf{Status / Note} \\
\hline
Shoulder push/pull strength & $r \approx 0.56$--$0.60$ & Mixed competitive & 4 & EVIDENCE\_BACKED; no multivariate partial-$r$ isolating shoulder from finger strength \\[4pt]
Hip flexibility (Lattice) & $R = 0.026$ ($p = 0.62$) & $N = 422$ bouldering & 5 & Null for grade; BIOPLAUSIBLE for technique facilitation \\[4pt]
Core strength (L-sit/front lever) & Not reported & --- & 0 & TO\_VERIFY; no peer-reviewed grade-correlation study published \\[4pt]
Lower body CMJ & Null (bouldering/sport) & Multiple studies & 4 & EVIDENCE\_BACKED null; significant negative predictor speed climbing ($r = -0.78$ to $-0.90$) \\[4pt]
Years climbing & Positive, major & Mermier 2000 ($n=44$) \cite{MJPS00}; 8a.nu OLS & 5 & EVIDENCE\_BACKED; training-component dominance ($58.9\,\%$ variance) \\[4pt]
Rate of Force Development (RFD) & ES $= 1.02$--$1.58$ elite vs.\ advanced & \cite{SVS+21} & 4 & EVIDENCE\_BACKED for group discrimination; BIOPLAUSIBLE as grade predictor (no prospective design) \\
\hline
\end{tabular}
\end{table}

\paragraph{Shoulder strength.}
Upper-limb pulling and pushing strength (assessed on a shoulder crank ergometer) correlates $r \approx 0.56$--$0.60$ with redpoint grade in competitive climbers (EVIDENCE\_BACKED, ESS 4).
No multivariate study has reported the partial correlation of shoulder strength with grade after controlling for finger strength; shoulder strength may share substantial variance with general training volume and MFS.

\paragraph{Hip flexibility.}
The Lattice Training (2025) multivariate dataset ($N = 422$ bouldering) yielded a flexibility coefficient $R = 0.026$ ($p = 0.62$), constituting an effectively null grade predictor.
Flexibility is mechanistically plausible for enabling high-feet technique and reducing energy expenditure on balance-dependent moves, but this effect is not captured in cross-sectional grade correlations.

\paragraph{Core strength.}
No peer-reviewed study has published a validated grade correlation for core strength measures (L-sit hold time, front lever, hanging leg raise).
This constitutes a notable gap given the frequency with which core strength features in climbing training recommendations.

\begin{toverifybox}{Core Strength as a Grade Predictor · \toverify · ESS 0}
No published cross-sectional or prospective study provides a Pearson or partial correlation between an objectively measured core strength index and climbing grade in trained climbers.
Mechanistic rationale for body tension efficiency is plausible but unquantified.
\end{toverifybox}

\paragraph{Lower body power (CMJ).}
Countermovement jump height does not correlate significantly with bouldering or sport climbing grade in available cross-sectional studies (EVIDENCE\_BACKED null, ESS 4); elite sport climbers do not exhibit exceptional lower-body power relative to the broader athletic population.
In speed climbing, leg power is a significant negative predictor of completion time ($r = -0.78$ to $-0.90$), reflecting the sprint-like demands of speed routes; speed and difficulty climbing should not be treated as equivalent in terms of physical determinants.

\paragraph{Rate of Force Development (RFD).}
Stien et al.\ (2021): elite climbers showed RFD effect sizes $1.02$--$1.58$ SDs above advanced climbers on finger-specific dynamometry \cite{SVS+21}, establishing RFD as a discriminant between ability groups (EVIDENCE\_BACKED, ESS 4).
Levernier and Laffaye (2019) demonstrated that 4 weeks of finger grip training increased RFD significantly in elite and top-ranked climbers ($n$ small; \cite{LL19}), providing limited intervention support.
RFD has not been measured as a prospective grade predictor in a longitudinal design; its predictive role beyond group discrimination is BIOPLAUSIBLE.

\paragraph{Years climbing.}
Training experience, operationalised as years climbing combined with strength metrics, constituted the dominant component in Mermier et al.\ (2000): the training-variable cluster explained $58.9\,\%$ of performance variance \cite{MJPS00}.
An OLS regression on the 8a.nu self-report database identified years climbing as a statistically significant positive predictor across all grade sub-groups.
Mechanistically, accumulated deliberate practice subsumes multiple adaptations (technique, connective tissue, MFS, CF) that are measured separately in cross-sectional profiling; years climbing is therefore a latent composite rather than an independent causal variable.

\subsubsection{Cross-Sectional Limitations: No Causal Inference}

\begin{evidencebox}{Cross-Sectional Limitations: No Causal Inference · \evidencebacked · ESS 7}
All predictor correlations reported in Sections~\ref{sec:cfpredictor} and~\ref{sec:anthropometrics} derive from cross-sectional designs.
Three systematic confounds apply:

\begin{enumerate}
  \item \textbf{Reverse causation}: elite athletes score higher on every measured capacity partly because higher training volume produces both elite grades and higher test scores simultaneously.
        High MFS, CF, and RFD in elite climbers does not imply that training these variables produces elite performance.
  \item \textbf{Range restriction}: studies recruiting ``elite'' or ``competitive'' sub-samples observe inflated correlations compared with grade-matched samples spanning a narrower performance band; correlations from mixed recreational--elite samples are not directly applicable to elite-only training decisions.
  \item \textbf{Grade self-report noise}: most studies rely on participant-reported redpoint grade, which conflates gym vs.\ outdoor ratings, redpoint vs.\ onsight performance, and regional grade inflation.
\end{enumerate}

Predictors with published \emph{intervention} evidence for the underlying physical quality (not climbing grade):
MFS increases have been demonstrated by Hermans et al.\ \cite{HO22} and meta-analytically by Stien et al.\ \cite{SO23};
RFD increases by Levernier \& Laffaye \cite{LL19}.
\textbf{Climbing grade was not the primary outcome in any of these trials}; transfer to on-wall performance remains an assumption supported by cross-sectional correlation, not by controlled longitudinal evidence.
\end{evidencebox}
%---------------------------------------------------------------------
\subsection{Multivariate Models and Probabilistic Grade Estimation}
\label{sec:multivariate}

\subsubsection{Best Available Multivariate Models}

Table~\ref{tab:multiv_full} summarises peer-reviewed multivariate models predicting climbing grade from physical and training data. No prospective (longitudinal) multivariate model predicting grade \emph{improvement} has been published; all entries are cross-sectional.

{\footnotesize
\setlength{\tabcolsep}{5pt}
\renewcommand{\arraystretch}{1.2}
\begin{longtable}{>{\raggedright\arraybackslash}p{3.0cm}
                  >{\raggedright\arraybackslash}p{1.3cm}
                  >{\raggedright\arraybackslash}p{2.4cm}
                  >{\raggedright\arraybackslash}p{1.4cm}
                  >{\raggedright\arraybackslash}p{6.0cm}}
\caption{Multivariate grade-prediction models: peer-reviewed studies.}\label{tab:multiv_full}\\
\toprule
\textbf{Study} & \textbf{$n$} & \textbf{Method} & \textbf{$R^{2}$} & \textbf{Notes} \\
\midrule
\endfirsthead
\multicolumn{5}{c}{\textit{(continued)}}\\
\toprule
\textbf{Study} & \textbf{$n$} & \textbf{Method} & \textbf{$R^{2}$} & \textbf{Notes} \\
\midrule
\endhead
\bottomrule
\endlastfoot

Mermier et al.\ \cite{MJPS00} & 44 mixed-sex (5.6--5.13c) & PCA + stepwise regression & 0.66 & Training component 58.9\%; flexibility 19.8\%; anthropometrics 9.8\%; finger strength not isolated --- training volume proxies cumulative finger loading \\[3pt]

Bala\v{s} et al.\ \cite{BPMC11} & $\sim$60 mixed-sex & Structural equation modelling (SEM) & $\sim$0.97$^{a}$ & Latent hand-arm strength + body fat + climbing experience; inflated $R^{2}$ plausible through SEM model specification --- see note $a$ \\[3pt]

Laffaye et al.\ 2016 \cite{TEMP_Laffaye2016} & 41 & PCA + regression & 0.64 & Training component (fingers + upper-limb power) 46\%; DOI unconfirmed (\toverify) \\[3pt]

Saeterbakken et al.\ \cite{TEMP_Saeterbakken2023} & 74--142 (by discipline) & Multiple regression & 0.58 & Women lead: upper-body power $\beta = 0.59$, finger endurance $\beta = 0.48$; men lead: lower-body + upper-body + fat + motor-planning $\beta = 0.27$--$0.36$; discipline-specific models \\[3pt]

Stien et al.\ \cite{BMK+25} & 307 (185M + 122F) & CFA & $\sim$0.96 & Finger strength 68\%/64\% + intermittent endurance 28\%/34\% (M/F); model includes only these two domains --- additional predictors not evaluated \\[3pt]

\end{longtable}
\noindent$^{a}$\textit{Bala\v{s} et al.\ \cite{BPMC11}: $R^{2} \approx 0.97$ is SEM-derived on $n \approx 60$. SEM can yield high fit indices through model specification choices that do not reflect true shared variance; the model has not been externally cross-validated. Treat as exploratory.}
}

The models collectively identify three consistently dominant predictor domains: (1) finger strength or cumulative finger-loading history; (2) upper-body pulling capacity; (3) body composition (mass, body fat). Training-volume variables frequently explain the largest single variance component (\cite{MJPS00}: 58.9\%) but are retrospective self-report and confound technique acquisition with physiological adaptation. No published model incorporates technique, movement efficiency, or route-reading as quantified variables.

\subsubsection{Machine-Learning and Citizen-Science Approaches}

Several non-peer-reviewed analyses apply supervised machine-learning to self-reported climber datasets.

\begin{folklorebox}{ML Grade Prediction: Citizen-Science Analyses · TO\_VERIFY · ESS 1}
\begin{itemize}
  \item Thomas C.\ Georgiou (Medium, 2025)\footnote{\url{https://medium.com/@tcgeorgiou/predicting-climbing-grades-with-machine-learning}} applied Ridge Regression, Gradient Boosting, and Random Forest to a self-reported climber dataset including max hang, repeater endurance, and pull-up max. Feature importance consistently ranked max hang and strength-to-weight first; body anthropometrics contributed minimally. No hold-out validation reported; dataset size and source not specified.
  \item Steve Bachmeier (GitHub, date unclear)\footnote{\url{https://github.com/stevebachmeier/climbing-grade-predictor}} fit OLS regression on the 8a.nu dataset ($\sim$400\,000 ascents). Years of climbing experience was the strongest grade predictor; height was a negative predictor for males (taller climbers averaged lower boulder grades). Confounding by experience with grade ceiling not controlled.
\end{itemize}
These analyses are citizen-science, non-peer-reviewed, and their datasets are self-selected with self-reported grades. Results are consistent with the peer-reviewed literature in ranking strength-to-weight first, but methodology does not permit causal inference.
\end{folklorebox}

\subsubsection{Unexplained Grade Variance}

\begin{bioplausbox}{Sources of Unexplained Grade Variance · BIOPLAUSIBLE · ESS 2}
Even the best physical-only multivariate models leave 33--50\% of grade variance unexplained ($1 - R^{2}$ from Table~\ref{tab:multiv_full}). Principal candidate sources of this residual:
\begin{itemize}
  \item \textbf{Technique and movement efficiency.} Footwork precision, weight transfer, body flagging, and rest-position identification are near-impossible to quantify without motion-capture instrumentation on standardised test routes.
  \item \textbf{Psychological performance.} Fear of falling, competition anxiety, self-efficacy, and sustained effort under discomfort are mechanistically linked to performance via arousal regulation and attentional focus, but no climbing study has quantified their independent grade contribution in a multivariate physical + psychological model (\toverify).
  \item \textbf{Route-reading and tactical decision-making.} Beta selection on unfamiliar terrain affects energy expenditure per move; no quantitative proxy is included in any published physical model.
  \item \textbf{Grade subjectivity and self-report noise.} Gym versus outdoor grade discrepancies, grade inflation across crags, and redpoint-versus-onsight differences introduce systematic and random error in the dependent variable of every reviewed model.
  \item \textbf{Training specificity mismatch.} Athletes trained predominantly on campus boards and hangboards accumulate high test scores but may have reduced climbing coordination relative to wall-trained athletes at equivalent physical capacity --- an artefact not separable from true ability in cross-sectional designs.
\end{itemize}
Quantifying these contributors requires technique-capture methodology (motion analysis, force-instrumented holds), validated psychological instruments administered concurrently with physical testing, and outdoor grade verification by a blinded assessor --- none of which have been achieved in a single published study.
\end{bioplausbox}

\subsubsection{Probabilistic Grade Estimation: IRT, Rasch Models, and Elo Systems}

\paragraph{Theoretical framework.}
Item Response Theory (IRT) and Rasch models place climber ability ($\theta$) and route difficulty ($\beta$) on the same interval-scale latent dimension. The 1PL (Rasch) model gives success probability $P(\text{send}) = \exp(\theta - \beta) / [1 + \exp(\theta - \beta)]$, mathematically equivalent to the Bradley--Terry paired-comparison model and to the Elo rating function used in chess. Applied to binary send/fail logbook data, these models simultaneously estimate $\theta$ for each climber and $\beta$ for each route without requiring a common rating committee --- a property well-suited to the decentralised structure of climbing grade data.

\paragraph{Whole-History Rating applied to climbing.}
Scarff (2020)\footnote{Scarff, H.\ ``Whole-History Rating for Rock Climbing,'' arXiv:2001.05388 (preprint, not peer-reviewed).} applied the Whole-History Rating (WHR) model --- a time-varying extension of the Bradley--Terry model --- to logbook ascent data. The model iteratively estimates both climber ability and route difficulty from binary send records. Scarff's analysis found that V-scale grade increments correspond to approximately $3.17\times$ the success probability odds ratio per grade step, providing an empirical difficulty spacing for the V-scale. This preprint has not undergone peer review; findings should be treated as exploratory.

\paragraph{Bayesian Route Response Theory.}
Rosenthal (2022)\singlesource\footnote{Rosenthal, D.\ ``Bayesian Rock Climbing Rankings,'' data science blog post, 2022. URL: \url{https://www.bestclimbing.com/bayesian-rankings} (verify URL independently).} applied a Bayesian Rasch model to 8a.nu ascent data, simultaneously estimating climber ability and route difficulty on a latent scale. The approach is structurally equivalent to Bayesian 1PL-IRT. This is a web-only source with no independent peer review; the analysis is cited for methodological description only.

\paragraph{Frontiers 2024: sensor-based grade classification.}
Michailov et al.\ \cite{OM25} used surface EMG of the forearm flexors and body-worn accelerometers to classify routes into three difficulty levels, reporting 98\% classification accuracy. Forearm flexor EMG amplitude exhibited a logarithmic relationship with route grade. The study operationalises grade as a classifier outcome rather than a continuous latent trait, and uses $n < 20$ participants on a controlled indoor wall --- not a true IRT/Rasch implementation, but the first peer-reviewed study to link objective physiological signals to grade classification using ML methods.

\paragraph{Bayesian MCMC on theCrag data.}
Drummond (2021)\singlesource\ fit a Bayesian MCMC model to theCrag/Vertical-Life ascent data and independently recovered a V-scale difficulty increment of $3.17\times$ per grade step --- consistent with Scarff (2020) despite using different data sources and model implementation. This result has not been published in a peer-reviewed venue.

\medskip
No full IRT or Rasch analysis specifically designed as a \emph{grade-prediction} tool for individual climbers has been published in a peer-reviewed journal. Current applications are exploratory data-science projects. The theoretical fit of Rasch/1PL models to climbing logbook data is strong (binary outcomes, no time pressure on the route-difficulty estimate), making formal IRT validation a tractable near-term research question.

\subsubsection{Big-Data 8a.nu and Moonboard Analyses}

\begin{folklorebox}{Big-Data Citizen-Science Grade Analyses · FOLKLORE\_VERIFIED}
Multiple independent analyses of the publicly scraped 8a.nu dataset ($\geq$400\,000 ascent records, Kaggle and GitHub) produce consistent findings across $\geq$3 independent sources:
\begin{itemize}
  \item OLS regression on 8a.nu data predicts grade with RMSE $\approx 4$ grade units (roughly one full Fontainebleau grade); years of climbing experience is the strongest predictor; BMI and height show minimal independent contribution after controlling for experience.
  \item Moonboard route grade prediction (Stanford CS229 2017 report)\footnote{Stanford CS229 project report, 2017. Citizen-science; not peer-reviewed.}: CNN achieved highest accuracy among compared methods but remained below human-expert accuracy. Features extracted from hold-type and position maps; grade labels are community-assigned.
  \item ``Board-to-Board'' (arXiv 2023)\footnote{arXiv preprint, 2023. Not peer-reviewed.}: CNN model trained on one Moonboard generation showed substantial accuracy drops when applied to a different Moonboard version, indicating that learned representations are board-version-specific rather than general climbing difficulty features.
\end{itemize}
All analyses are non-peer-reviewed, use self-reported or community-assigned grades as ground truth, and apply to structured board climbing (Moonboard) or self-selected outdoor/indoor ascents (8a.nu). Generalisability to competition climbing or individual grade prediction is unconfirmed.
\end{folklorebox}

%---------------------------------------------------------------------
\subsection{Test Batteries in Research and Practice}
\label{sec:testbatteries}

\subsubsection{Overview and Comparison}

Table~\ref{tab:batteries} compares the principal test batteries used in research and applied settings. No battery has been validated prospectively (predicting grade change over a training cycle); all validity data are cross-sectional or within-session reliability.

{\footnotesize
\setlength{\tabcolsep}{5pt}
\renewcommand{\arraystretch}{1.2}
\begin{longtable}{>{\raggedright\arraybackslash}p{2.8cm}
                  >{\raggedright\arraybackslash}p{3.8cm}
                  >{\raggedright\arraybackslash}p{2.2cm}
                  >{\raggedright\arraybackslash}p{2.5cm}
                  >{\raggedright\arraybackslash}p{2.5cm}}
\caption{Test batteries in climbing research and practice.}\label{tab:batteries}\\
\toprule
\textbf{Battery} & \textbf{Components} & \textbf{Population} & \textbf{Published validity} & \textbf{Norms available} \\
\midrule
\endfirsthead
\multicolumn{5}{c}{\textit{(continued)}}\\
\toprule
\textbf{Battery} & \textbf{Components} & \textbf{Population} & \textbf{Published validity} & \textbf{Norms available} \\
\midrule
\endhead
\bottomrule
\endlastfoot

Lattice athlete profiling & 7\,s 2-arm hang 20\,mm \%BW; 2RM weighted pull-up \%BW; 60\% MFS power endurance TTE (7/3); box-split flexibility & Self-selected, all levels & Blog-level regression ($n = 422$ routes, $n = 901$ bouldering) & Internal; community-reverse-engineered \\[3pt]

9c test & 5\,s max hang 20\,mm \%BW; 1RM weighted pull-up \%BW; core tier (bent-knee L-sit to front lever); bar dead-hang time & Community, any level & None published & Community calculators only \\[3pt]

Stien 2022 mini-review battery \cite{SSA22} & 7\,s max hang; continuous hang to failure (20--25\,mm); 7/3 60\% MVC to failure; isometric pull-up for RFD on 23\,mm campus rung & Research & Recommended protocol only & None \\[3pt]

Climbro 4-test & MFS; explosive RFD pull; continuous TTE; intermittent 7/3 TTE & Research + field & Validated across 3 IRCRA performance groups\footnote{Climbro commercial system. Web: \url{https://climbro.com} (single-source; independent peer-reviewed validation pending).} & Emerging \\[3pt]

\end{longtable}
}

\subsubsection{The 9c Test: Protocol and Evidence Evaluation}

\begin{folklorebox}{9c Test: Protocol · FOLKLORE\_VERIFIED · ESS 0}
The 9c test is a community-originated 4-component battery mapping a 4--40 point composite score to a predicted climbing grade. Protocol verified across $\geq$3 independent sources (Hooper's Beta YouTube analysis, UKC forums, theCrag community discussions, Men's Health coverage). Each component scores 1--10 points; total score determines the predicted grade.

\medskip
\noindent\textbf{Exercise 1 --- Maximum finger strength} (5\,s weighted hang, 20\,mm edge, half-crimp, added weight as \%BW):
\smallskip

\begin{tabular}{llllllllll}
1 pt & 2 & 3 & 4 & 5 & 6 & 7 & 8 & 9 & 10 pt \\
100\% & 110\% & 120\% & 130\% & 140\% & 150\% & 160\% & 180\% & 200\% & 220\%BW \\
\end{tabular}

\medskip
\noindent\textbf{Exercise 2 --- Maximum pull-up} (1RM weighted, expressed as \%BW):
\smallskip

\begin{tabular}{llllllllll}
1 pt & 2 & 3 & 4 & 5 & 6 & 7 & 8 & 9 & 10 pt \\
100\% & 110\% & 120\% & 130\% & 140\% & 150\% & 160\% & 180\% & 200\% & 220\%BW \\
\end{tabular}

\medskip
\noindent\textbf{Exercise 3 --- Core (tiered, ascending):}
\smallskip

\begin{tabular}{>{\raggedright}p{2.0cm} >{\raggedright\arraybackslash}p{10cm}}
1 pt & 10\,s bent-knee L-sit \\
4 pt & 10\,s straight L-sit \\
7 pt & 5\,s front lever \\
10 pt & 30\,s front lever \\
\end{tabular}

\medskip
\noindent\textbf{Exercise 4 --- Bar dead-hang (unweighted):}
\smallskip

\begin{tabular}{ll}
1 pt = 10\,s & 10 pt = 100\,s \\
\end{tabular}

\medskip
\noindent\textbf{Grade prediction from total score (1--40):}
\smallskip

\begin{tabular}{ll}
\toprule
\textbf{Score} & \textbf{Predicted grade (Sport, French)} \\
\midrule
40 & 9c \\
36--39 & 9a+/9b+ \\
32--35 & 9a/9a+ \\
28--31 & 8c/8c+ \\
24--27 & 8b/8b+ \\
20--23 & 8a/8a+ \\
16--17 & 7c \\
12--15 & 7b/7b+ \\
\bottomrule
\end{tabular}
\end{folklorebox}

\paragraph{Evidence evaluation of individual components.}
Exercise~1 (finger strength, 5\,s hang): \evidencebacked\ as a grade predictor --- \cite{BBW25} reports $r = 0.89$ with bouldering grade for a climbing-specific isometric hang. Exercise~2 (weighted pull-up): \evidencebacked\ as a moderate grade predictor --- \cite{BBW25} reports $r = 0.55$ for bouldering. Exercise~3 (core/front lever): \toverify --- no peer-reviewed study has measured grade correlation with front-lever performance or timed L-sit; the inclusion is \bioplausible\ on the mechanistic basis that anterior-chain stiffness reduces energy expenditure during body-tension movements on overhanging terrain, but no controlled climbing study has quantified this link. Exercise~4 (bar dead-hang): \toverify\ --- tests grip endurance on a large jug in a fully extended arm position; this is mechanistically dissimilar from the finger-flexor fatigue pattern during climbing (half-crimp or open-hand on edges); no published study has correlated bar dead-hang time with climbing grade.

No peer-reviewed study has validated the 9c composite score against climbing grade. This absence was confirmed by a comprehensive search of PubMed, IJSPP, and Frontiers databases. ESS~0 applies to the composite battery as a grade-prediction instrument.

\paragraph{Injury-risk note.}
The finger-strength scoring increments jump 20\%BW per point at the upper scale (160\% $\to$ 180\% $\to$ 200\% $\to$ 220\%BW). Hooper's Beta\singlesource\ notes that pursuit of these load targets as training goals poses meaningful pulley injury risk; the jump from 160\% to 180\%BW represents a 12.5\% relative force increase on the A2 pulley at a loading level already at or above the estimated A2 failure threshold for some athletes. These values should not be used as training targets without progressive load management (risk stratum: \textbf{High} $>$20\% relative load increase; see Section~\ref{sec:injury}).

\paragraph{Predictive range limitations.}
The 9c test's discrimination is coarsest at the extremes. Below $\sim$7a (French), technique dominates over measurable physical attributes, and physical scores in the 12--19 point range offer little actionable precision. Above 8c+ ($\sim$score 36+), only 3--4 points separate adjacent grade bands, placing adjacent grade predictions within measurement noise for all four components.

\paragraph{What the 9c test does not measure.}
The four components contain no assessment of forearm aerobic endurance (critical force / W$'$), technique, footwork, route-reading, psychological performance, flexibility, or skin condition. Power endurance (the capacity to sustain near-maximal effort through a sequence of hard moves) is also absent --- a domain directly predictive of sport-climbing redpoint performance.

\subsubsection{Comparison: 9c Test vs.\ Lattice Profiling}

\begin{evidencebox}{9c Test vs.\ Lattice Battery: Structural Comparison · EVIDENCE\_BACKED · ESS 2}
\begin{tabular}{>{\raggedright}p{4.0cm}>{\raggedright}p{4.0cm}>{\raggedright\arraybackslash}p{4.5cm}}
\toprule
\textbf{Dimension} & \textbf{9c test} & \textbf{Lattice profiling} \\
\midrule
Finger strength & Ex.\ 1 (5\,s, 20\,mm, \%BW) & 7\,s 2-arm hang (20\,mm, \%BW) \\
Pull-up strength & Ex.\ 2 (1RM, \%BW) & 2RM \%BW \\
Power endurance & Absent & 60\% MFS TTE (7/3 intermittent) \\
Flexibility & Absent & Box split \\
Core / front lever & Ex.\ 3 (tiered) & Absent \\
Bar dead-hang & Ex.\ 4 & Absent \\
Grade output & Composite score mapped to grade & Profiling vs.\ normative percentiles \\
Normative validation & None published & Blog-level regression ($n > 400$) \\
Purpose & Classificatory & Diagnostic + prescriptive \\
\bottomrule
\end{tabular}

\medskip
\noindent Lattice uniquely quantifies power endurance, the strongest endurance-domain predictor in route models \cite{TEMP_Saeterbakken2023}. The 9c test uniquely includes bar dead-hang endurance and core tiering, neither of which has published grade-correlation evidence. Lattice battery has blog-level regression validation against grade and normative data from $>$400 climbers; 9c composite has zero published validation.
\end{evidencebox}

\subsubsection{Lattice Training Athlete Profiling}

Lattice Training's applied battery comprises four components assessed under standardised conditions: (1) 7\,s 2-arm dead-hang on a 20\,mm edge at \%BW (MFS proxy); (2) 2-repetition-maximum weighted pull-up at \%BW; (3) 60\% MFS power endurance --- intermittent 7\,s-on / 3\,s-off protocol to volitional failure (TTE); (4) box-split flexibility. Normative data are published at blog level for bouldering ($n = 901$) and routes ($n = 422$) as of 2025\singlesource\ (\url{https://latticetraining.com/2-arm-fs-test/}). No independent peer-reviewed validation of the Lattice normative dataset has been published. Community reverse-engineering of Lattice norms (WMG Climbing blog\singlesource) corroborates the \%BW values published directly by Lattice.

The Lattice framework is explicitly prescriptive: test output is compared against percentile norms by grade band and discipline, identifying which physical domain is most limiting relative to grade-matched peers. This population-referenced interpretive model is absent from the 9c test.

\subsubsection{Research-Grade Batteries}

Stien et al.\ \cite{SSA22} (mini-review, $n$ = synthesis, not original data) recommend a 4-test research battery: (1) 7\,s maximum isometric hang on a 20--23\,mm edge; (2) continuous hang to failure on the same edge; (3) 7\,s-on / 3\,s-off at 60\% MVC to volitional failure; (4) isometric pull-up for RFD on a 23\,mm campus rung at $\geq$200\,Hz sampling. Component (4) targets RFD --- the discriminating variable between elite and advanced climbers identified by Stien et al.\ \cite{SVS+21} (ES = 1.02--1.58) --- and requires research-grade force platforms, not consumer devices (see RFD measurement section above).

Langer et al.\ \cite{LSW23} (systematic review, $k$ = studies synthesised) found that climbing-specific finger-strength tests outperform non-specific handgrip dynamometry in multivariate grade-prediction models, supporting sport-specific over general testing. Reliability coefficients (ICC) were highest for 20--23\,mm edge protocols.

Giles et al.\ \cite{GRT06} provide a foundational narrative review of climbing physiology that remains a methodological reference for test-variable selection, particularly for aerobic and anaerobic endurance components. The Giles 2019 protocol \cite{GCT+19} operationalises critical force (CF) and anaerobic work capacity ($W'$) estimation from a two-bout sustained contraction protocol on a 20\,mm edge, enabling field estimation of aerobic finger-flexor capacity without laboratory equipment.

\subsubsection{Community Tools}

\begin{toverifybox}{Community Grade-Prediction Tools · TO\_VERIFY · ESS 0--1}
The following tools are available for community use. None has undergone independent peer-reviewed predictive validity testing.
\begin{itemize}
  \item \textbf{StrengthClimbing Finger Strength Analyzer 2.0} (J.\ Banaszczyk)\footnote{\url{https://strengthclimbing.com/finger-strength-analyzer/}}: inputs 20\,mm max hang and weighted pull-up; outputs predicted boulder grade. Methodology blog-described only; no hold-out validation published.
  \item \textbf{StrengthClimbing Critical Force Calculator}\footnote{\url{https://strengthclimbing.com/critical-force-calculator/}}: implements the Giles et al.\ \cite{GCT+19} two-bout CF / $W'$ estimation equations. Calculation is mechanistically grounded \cite{GCT+19}; field protocol validity depends on accurate self-measurement.
  \item \textbf{OpenClimbing 9c calculator}: web implementation of the 9c scoring table above. No novel validation.
\end{itemize}
\end{toverifybox}

\subsubsection{Open Questions the Literature Cannot Yet Answer}

The following gaps were identified as unresolved across the reviewed literature:

\begin{enumerate}[leftmargin=1.5em]
  \item Does any field battery (9c test, Lattice profiling) predict grade \emph{improvement} over time? All published validity data are cross-sectional; no prospective battery-to-improvement study has been published.
  \item How does the relative importance of predictors (finger strength, pull-up, critical force) shift across the grade spectrum from V4 to V17? Current models aggregate across ability levels or report by discipline only; within-grade-band predictor weighting is unstudied.
  \item What is the independent contribution of critical force to grade prediction after simultaneous control for MFS in a single multivariate model? Current evidence uses sequential or stepwise entry; partial regression coefficients for CF vs.\ MFS in the same model are unavailable.
  \item Can technique and movement efficiency be reliably quantified in a field setting, and what fraction of the 33--50\% unexplained variance do they account for?
  \item Do predictor profiles differ by sex and discipline (bouldering vs.\ sport climbing) at elite competition level ($\geq$9a / V16)? Current studies are underpowered at elite level due to sparse sampling.
  \item Is the MFS--grade relationship linear across the full grade range, or does it exhibit a threshold or diminishing-returns structure above $\sim$V12 (8b)? Existing data span wide grade ranges without modelling non-linearity.
  \item Do structural anatomical variables (A2 pulley cross-sectional area, tendon insertion geometry, finger segment length) impose performance ceilings that training-induced strength gains cannot overcome? No intervention study has measured anatomical structure alongside MFS change.
\end{enumerate}
% -------------------------------------------------------------------
